%% ============================================================
%% 分片 31：原文第 11 节「$A$-Branes and ${\cal D}$-Modules」（原文 glshort.tex 起于第 9938 行）
%% 由 tools/harvmac2tex.py 从 arXiv 源码（hep-th/0604151v3）自动生成。
%% 只译英文散文；公式、编号、\label、\cite 一律不改。
%% ============================================================

\section{$A$-Branes and ${\cal D}$-Modules}
\label{sec:abranes}

 
  

  
  

  

  

A crucial ingredient is missing from what we have presented so
far. According to the geometric Langlands program, there is a
natural correspondence between homomorphisms of $\pi_1(C)$ into
$^L\negthinspace G$ and Hecke eigensheaves on ${\EUM}(G,C)$ which
are also ${\cal D}$-modules.  A ${\cal D}$-module is a module for
the sheaf of differential operators; the concept is further
elucidated below.

Concerning the left hand side of the correspondence, we learned in
section 第~\ref{sec:electricbranes}~节\ that homomorphisms of $\pi_1(C)$ to
$^L\negthinspace G$ correspond to electric eigenbranes.  Their
duals, which are branes of type ${\CMF}$ -- that is, branes
wrapped on a fiber of the Hitchin fibration, with a flat
Chan-Paton bundle -- must therefore be magnetic eigenbranes.
This statement follows from $S$-duality and can also be verified
directly \cite{ref26}, using  the abelianization that is provided by
the spectral cover construction.

However, the magnetic eigenbranes that we get this way are branes
of type $(A,B,A)$ on $\MH(G,C)$, while in the conventional
statement of the geometric Langlands correspondence, the right
hand side is supposed to involve instead a ${\cal D}$-module on,
roughly speaking, $\M(G,C)$, the moduli space of $G$-bundles.

This section is devoted to supplying the missing link. We will
show that every $A$-brane in complex structure $K$ on $\MH(G,C)$
automatically gives rise to, roughly speaking, a ${\cal D}$-module
on $\M(G,C)$.

One can regard ${\cal D}$-modules on $\M(G,C)$ as $B$-branes on a
noncommutative deformation of the cotangent bundle of $\M(G,C)$.
It has been argued previously \cite{ref41}\ that for a certain class of
holomorphic symplectic manifolds, the category of $A$-branes is
equivalent to the category of $B$-branes on the noncommutative
deformation of the same manifold. The connection between
$A$-branes on $\M_H(G,C)$ and $\cal D$-modules on $\M(G,C)$ can be
viewed as a special case of this.

Actually, in the geometric Langlands program, it is not sufficient
(nor is it possible, because of the way the Hecke operators act)
to work only on $\M(G,C)$, the moduli space of stable bundles. One
instead constructs a ${\cal D}$-module on ${\rm Bun}_G(C)$, the
``stack'' of all $G$-bundles on $C$.  The gauge theory analog of
the stack of all $G$-bundles is simply the space of all
connections. The analog in our construction of the usual statement
that the ${\cal D}$-modules are defined on ${\rm Bun}_G(C)$ is
that the branes of interest can be defined in the underlying gauge
theory, and not only in the sigma-model with target space
$\MH(G,C)$. Showing this for the important branes is in fact the
goal of section 第~\ref{sec:zapendix}~节.
